\chapter{Residual correction of VaR and ES}
\label{ch:recalibration}

Widening a predictive distribution does not determine the probability of a severe loss, even when its scale follows the turbulent conditions illustrated earlier. We retain the issued scale and estimate its tail conversion from past outcomes. The procedure below is evaluated in Chapter~\ref{ch:tailresults}.

\section{Why correct the residual tail?}

At level $\alpha$, Value-at-Risk (VaR) defines a return threshold and Expected Shortfall (ES) averages the lowest $\alpha$ fraction of outcomes. Negative returns represent losses. A forecast may follow changing scale while assigning insufficient probability to severe losses, so variance accuracy alone cannot establish accurate VaR and ES.

We take the magnitude of the issued negative VaR as a positive scale and divide each past return by the scale issued before it occurred. Here \emph{residuals} means these standardised outcomes, rather than mean-subtracted regression errors. Their empirical lower quantile and tail average supply the replacement conversion:

\begin{equation}
  \text{new tail pair}=\text{current scale}\times\text{fitted residual tail pair}.
  \label{eq:recal-workflow}
\end{equation}

For example, residual coefficients $-1.2$ and $-1.8$ applied to an original VaR magnitude of $2\%$ give replacement VaR $-2.4\%$ and ES $-3.6\%$. The same coefficients produce smaller losses when multiplied by a quieter day's smaller scale. After warm-up, original ES no longer independently determines the pair. The procedure therefore requires VaR rather than a complete density or standard deviation, although its success still depends on the information retained by that changing scale.

\begin{figure}[!htbp]
\centering
\begin{tikzpicture}[node distance=7mm and 8mm,box/.style={draw=black!45,rounded corners,align=center,text width=3.8cm,minimum height=14mm,font=\small},arr/.style={-{Latex[length=2mm]},thick,draw=black!60}]
\node[box,fill=blue!5] (old) {Past returns and their\\previously issued scales};
\node[box,right=of old,fill=blue!5] (res) {Divide each return by\\its own forecast scale};
\node[box,below=of res,fill=green!5] (tail) {Estimate the residual\\quantile and tail average};
\node[box,left=of tail,fill=green!5] (new) {Multiply both by\\today's forecast scale};
\draw[arr] (old)--(res);\draw[arr] (res)--(tail);\draw[arr] (tail)--(new);
\end{tikzpicture}
\caption[How residual tail correction works]{The correction retains a changing daily scale and estimates the tail from past standardised outcomes. Each scale was issued before its associated return. Forecast ordering and the score's numerical floor are specified in the text.}
\label{fig:tail-correction-workflow}
\end{figure}

With fixed tail shape, this is FHS expressed through VaR \citep{baroneadesi1999var,baroneadesi2008garch}. VaR's fixed multiple of volatility cancels when residual tails are restored to current scale. The forecasts coincide if histories, weights, quantile conventions and safeguards agree. Changing location or shape can break that relationship. Rescaling empirical tails is established; we evaluate its predictive benefit and limits.

\section{How the correction uses past observations}

We retain the original forecast for the first 250 observations. Thereafter, we estimate the residual quantile and integrated tail average before issuing each next pair. After observing the return, we divide it by the original scale issued for that date and add the resulting residual to the history. Keeping those scales preserves the information available at each forecast date. We do not restandardise an earlier return using today's scale, since the residual must measure the outcome relative to the forecast that preceded it.

The integrated average uses exactly the requested sample fraction. For a 1\% tail with 250 residuals, it averages the two lowest values and half the third, dividing by 2.5. Fractional boundary weights prevent excess mass from ties or non-integer tail counts, but create no new observations. Even 500 residuals provide only five observations for a 1\% tail, illustrating the limited severity information at rarer levels.

Numerical safeguards keep ES below VaR and scores finite near zero. These constraints form part of the evaluated procedure; saved losses alone cannot certify that they never activate. Appendix~\ref{sec:tail-score-details} gives their values and the original-model safeguards, and identifies retrospectively fitted, unscored initial histories in some market procedures.

\section{Choosing a history and deciding whether to correct}

Residual-history length controls an estimation trade-off. Full history retains every available residual after initialisation; a rolling window removes older values at a fixed length. More observations may improve precision under stability, but outdated residuals can describe the wrong tail after a change. Shorter windows can adapt sooner while containing fewer extremes. We examine these consequences without assuming an optimal length.

The fixed market study compares full history with 1,000- and 2,000-residual windows. Its \emph{guard} selects correction only when recent mean loss is sufficiently lower than the original forecast's. After 250 post-warm-up comparisons, it checks up to 1,000 past scores against fixed tail-specific margins (Appendix~\ref{sec:tail-score-details}). Both candidates continue updating when unselected, so the next decision can compare their recent performance even if one has not been issued. The guard selects whether to issue the fitted pair; it does not change the residual estimator. Previous performance informs the next choice without guaranteeing improvement.

Development selects the 1,000-residual correction for the later fixed guard. The exploratory selector considers several histories and was designed after evaluation-period exposure. Its additional choices answer a different question and cannot supply independent confirmation. Appendix~\ref{sec:tail-score-details} retains their candidate sets.

\section{What counts as an improvement?}

A breach count measures how often returns fall below VaR, but not how severe those losses are. We therefore assess VaR and ES jointly with FZ0 \citep{fissler2016higher,patton2019dynamic}. For return $y$, threshold $q$, ES $e<0$ and level $\alpha$, the loss is

\begin{equation}
 S_\alpha(y,q,e)
 =-\frac{(q-y)^+}{\alpha e}+\frac{q}{e}+\log(-e)-1,
 \qquad x^+=\max(x,0).
 \label{eq:recal-score}
\end{equation}

We prefer lower FZ0 and interpret differences as loss units, rather than percentages of financial losses avoided. The recorded numerical floor is applied before scoring, with return units fixed within each comparison.

The simulations isolate two consequences of fitting a replacement. Correcting the known true pair measures the cost of estimation when no persistent tail error needs repair. Correcting Gaussian forecasts on heavy-tailed returns measures repair opportunity. Varying path length and changing the law then show how available history and stability affect both. The explanation applies to these settings without supplying a general performance bound.

The grouped comparison examines whether retaining daily scale detail matters. It replaces varying forecasts with one raw-return pair in each scale group, leaving fewer observations to estimate each group's tail. We distinguish online estimates from hindsight fits that use scored outcomes. The comparison follows motivation from \citet{foster2023calibeating} without replicating a general calibration guarantee.

The next chapter measures the benefit of repairing a restrictive tail and the cost of correcting an already accurate forecast. Joint tail scores assess both effects; separate coverage diagnostics assess breach frequency.